\newcommand{\accuracyTable}{{
\begin{table*}[t]
\centering
\captionof{table}{Performance of LLaMA-7B compressed by \sysname, \sysname \texttt{(W)}, and baselines under different compression ratio (corresponding weight memory is listed inside bracket) on two language modeling datasets (measured by perplexity  (\textcolor{mygreen}{$\downarrow$})), eight common sense reasoning datasets (six measured by both individual and average accuracy (\textcolor{mygreen}{$\uparrow$}), TruthfulQA measured by BLEU score (\textcolor{mygreen}{$\uparrow$}), and GSM8K measured by Exact Match Accuracy (\textcolor{mygreen}{$\uparrow$})). The best performance is marked in bold. The relative performance gain compared to the best-performing baseline is marked in green inside bracket.
}
\vspace{-2mm}
\resizebox{1\textwidth}{!}{%}
\begin{tabular}{ccccccccccccc}
\midrule
\multicolumn{1}{c|}{\textsc{Ratio (Mem.)}}       & \multicolumn{1}{c|}{\textsc{Method}}    & WikiText-2\textcolor{mygreen}{$\downarrow$} & \multicolumn{1}{c|}{C4\textcolor{mygreen}{$\downarrow$}} & Openb. & ARC\_e & WinoG. & HellaS.  & PIQA & MathQA & \textbf{Average\textcolor{mygreen}{$\uparrow$}}  & TruthfulQA\textcolor{mygreen}{$\uparrow$} & GSM8K\textcolor{mygreen}{$\uparrow$}      \\ \midrule
\multicolumn{1}{c|}{\color[HTML]{9B9B9B}0\% (13.5 GB)}  & \multicolumn{1}{c|}{\color[HTML]{9B9B9B}Original}  & {\color[HTML]{9B9B9B}5.68}     & \multicolumn{1}{c|}{\color[HTML]{9B9B9B}7.34}      & {\color[HTML]{9B9B9B} 0.34} & {\color[HTML]{9B9B9B} 0.75} & {\color[HTML]{9B9B9B} 0.70} & {\color[HTML]{9B9B9B} 0.57} & {\color[HTML]{9B9B9B} 0.79} & {\color[HTML]{9B9B9B} 0.27} & {\color[HTML]{9B9B9B} 0.57}    & {\color[HTML]{9B9B9B}0.30}  & {\color[HTML]{9B9B9B}0.09}\\ \midrule
\multicolumn{1}{c|}{\multirow{4}{*}{20\% (10.2 GB)}} & \multicolumn{1}{c|}{SVD}      & 20061                     & \multicolumn{1}{c|}{18800}          & 0.05  & 0.04  & 0.01  & 0.03  & 0.02  & 0.03  & 0.03  & 0.00   & 0.00 \\
\multicolumn{1}{c|}{}                     & \multicolumn{1}{c|}{FWSVD}     & 1727                     & \multicolumn{1}{c|}{1511}                      & 0.09  & 0.11  & 0.05  & 0.08  & 0.10  & 0.05  & 0.08  & 0.00   & 0.00 \\
\multicolumn{1}{c|}{}                     & \multicolumn{1}{c|}{ASVD}      & 11.14                   & \multicolumn{1}{c|}{15.93}                      & 0.29  & 0.53  & 0.64  & 0.41  & 0.68  & 0.17  & 0.45  & 0.21   & 0.04\\ \cmidrule{2-13} 
\multicolumn{1}{c|}{}                     & \multicolumn{1}{c|}{\xin{\sysname \texttt{(W)}}}  & 7.94 (\textcolor{mygreen}{$\downarrow$29\%})  & \multicolumn{1}{c|}{15.84 (\textcolor{mygreen}{$\downarrow$1\%})}       & 0.31  & 0.62  & 0.61 & 0.45  & 0.71  & 0.21  & 0.49 (\textcolor{mygreen}{$\uparrow$9\%})  & 0.26 (\textcolor{mygreen}{+0.05})  & 0.05 (\textcolor{mygreen}{+0.01})\\ 
\multicolumn{1}{c|}{}                     & \multicolumn{1}{c|}{\sysname}  & \textbf{7.73} (\textcolor{mygreen}{$\downarrow$31\%})  & \multicolumn{1}{c|}{\textbf{12.23} (\textcolor{mygreen}{$\downarrow$23\%})}       & \textbf{0.33}  & \textbf{0.67}  & \textbf{0.69 } & \textbf{0.55}  & \textbf{0.79}  & \textbf{0.26}  & \textbf{0.55} (\textcolor{mygreen}{$\uparrow$22\%})  & \textbf{0.28} (\textcolor{mygreen}{+0.07})  & \textbf{0.08} (\textcolor{mygreen}{+0.04})\\ \midrule

\multicolumn{1}{c|}{\multirow{4}{*}{40\% (7.76 GB)}} & \multicolumn{1}{c|}{SVD}      & 52489                    & \multicolumn{1}{c|}{47774}           & 0.04  & 0.04  & 0.05  & 0.01  & 0.03  & 0.02  & 0.03  & 0.00   & 0.00\\
\multicolumn{1}{c|}{}                     & \multicolumn{1}{c|}{FWSVD}     & 18156                     & \multicolumn{1}{c|}{12847}                    & 0.06  & 0.05  & 0.02  & 0.00  & 0.05  & 0.03  & 0.04  & 0.00   & 0.00\\
\multicolumn{1}{c|}{}                     & \multicolumn{1}{c|}{ASVD}      & 1407                 & \multicolumn{1}{c|}{1109}                          & 0.08  & 0.11  & 0.09  & 0.08  & 0.13  & 0.08  & 0.10  & 0.01   & 0.00\\ \cmidrule{2-13} 
\multicolumn{1}{c|}{}                     & \multicolumn{1}{c|}{\xin{\sysname \texttt{(W)}}}  & 13.73 (\textcolor{mygreen}{$\downarrow$99\%}) & \multicolumn{1}{c|}{75.42 (\textcolor{mygreen}{$\downarrow$93\%})}        & 0.25  & 0.33  & 0.55  & 0.40  & 0.63  & 0.12  & 0.38 (\textcolor{mygreen}{$\uparrow$280\%})  & 0.17 (\textcolor{mygreen}{+0.17})   & 0.02 (\textcolor{mygreen}{+0.02})\\ 
\multicolumn{1}{c|}{}                     & \multicolumn{1}{c|}{\sysname}  & \textbf{9.27} (\textcolor{mygreen}{$\downarrow$99\%}) & \multicolumn{1}{c|}{\textbf{15.63} (\textcolor{mygreen}{$\downarrow$99\%})}        & \textbf{0.29}  & \textbf{0.59}  & \textbf{0.68}  & \textbf{0.52}  & \textbf{0.69}  & \textbf{0.20}  & \textbf{0.50} (\textcolor{mygreen}{$\uparrow$400\%})  & \textbf{0.24} (\textcolor{mygreen}{+0.23})   & \textbf{0.07} (\textcolor{mygreen}{+0.07})\\ \midrule
\multicolumn{1}{c|}{\multirow{4}{*}{60\% (5.35 GB)}} & \multicolumn{1}{c|}{SVD}      & 105474                     & \multicolumn{1}{c|}{106976}        & 0.01  & 0.03  & 0.01  & 0.00  & 0.01  & 0.02  & 0.01  & 0.00   & 0.00\\
\multicolumn{1}{c|}{}                     & \multicolumn{1}{c|}{FWSVD}     & 32194                     & \multicolumn{1}{c|}{29292}                    & 0.06  & 0.02  & 0.01  & 0.01  & 0.02  & 0.03  & 0.03  & 0.00   & 0.00\\
\multicolumn{1}{c|}{}                     & \multicolumn{1}{c|}{ASVD}      & 57057                  & \multicolumn{1}{c|}{43036}                       & 0.05  & 0.04  & 0.06  & 0.09  & 0.08  & 0.05  & 0.06  & 0.00   & 0.00\\ \cmidrule{2-13} 
\multicolumn{1}{c|}{}                     & \multicolumn{1}{c|}{\xin{\sysname \texttt{(W)}}}  & 66.62 (\textcolor{mygreen}{$\downarrow$99\%}) & \multicolumn{1}{c|}{471.83 (\textcolor{mygreen}{$\downarrow$99\%})}       & 0.10  & 0.05  & 0.17  & 0.10  & 0.21  & 0.04  & 0.11 (\textcolor{mygreen}{$\uparrow$83\%})  & 0.01 (\textcolor{mygreen}{+0.01})   & 0.00 (\textcolor{mygreen}{+0.00})\\ 
\multicolumn{1}{c|}{}                     & \multicolumn{1}{c|}{\sysname}  & \textbf{15.00} (\textcolor{mygreen}{$\downarrow$99\%}) & \multicolumn{1}{c|}{\textbf{26.26} (\textcolor{mygreen}{$\downarrow$99\%})}       & \textbf{0.18}  & \textbf{0.42}  & \textbf{0.44}  & \textbf{0.31}  & \textbf{0.35}  & \textbf{0.12}  & \textbf{0.30} (\textcolor{mygreen}{$\uparrow$400\%})  & \textbf{0.14} (\textcolor{mygreen}{+0.14})   & \textbf{0.04} (\textcolor{mygreen}{+0.04})\\ \midrule
\multicolumn{1}{c|}{\multirow{4}{*}{80\% (2.58 GB)}} & \multicolumn{1}{c|}{SVD}      & 687291                     & \multicolumn{1}{c|}{708243}        & 0.00  & 0.01  & 0.02  & 0.01  & 0.01  & 0.00  & 0.01  & 0.00   & 0.00\\
\multicolumn{1}{c|}{}                     & \multicolumn{1}{c|}{FWSVD}     & 96872                    & \multicolumn{1}{c|}{89243}                     & 0.01  & 0.02  & 0.00  & 0.01  & 0.01  & 0.00  & 0.01  & 0.00   & 0.00\\
\multicolumn{1}{c|}{}                     & \multicolumn{1}{c|}{ASVD}      & 80425                 & \multicolumn{1}{c|}{67927}                        & 0.04  & 0.03  & 0.03  & 0.02  & 0.01  & 0.01  & 0.03  & 0.00   & 0.00\\ \cmidrule{2-13} 
\multicolumn{1}{c|}{}                     & \multicolumn{1}{c|}{\xin{\sysname \texttt{(W)}}}  & 1349 (\textcolor{mygreen}{$\downarrow$98\%}) & \multicolumn{1}{c|}{6224 (\textcolor{mygreen}{$\downarrow$91\%})}       & 0.07  & 0.03  & 0.04  & 0.02  & 0.07  & 0.01  & 0.04 (\textcolor{mygreen}{$\uparrow$33\%})  & 0.00 (\textcolor{mygreen}{+0.00})  & 0.00 (\textcolor{mygreen}{+0.00})\\ 
\multicolumn{1}{c|}{}                     & \multicolumn{1}{c|}{\sysname}  & \textbf{31.79} (\textcolor{mygreen}{$\downarrow$99\%}) & \multicolumn{1}{c|}{\textbf{43.71} (\textcolor{mygreen}{$\downarrow$99\%})}       & \textbf{0.11}  & \textbf{0.23}  & \textbf{0.21}  & \textbf{0.14}  & \textbf{0.17}  & \textbf{0.08}  & \textbf{0.16} (\textcolor{mygreen}{$\uparrow$433\%})  & \textbf{0.04} (\textcolor{mygreen}{+0.04})  & \textbf{0.02} (\textcolor{mygreen}{+0.02})\\ \midrule
\end{tabular}
}
\vspace{-4mm}
\label{tab:dataset_acc}
\end{table*}
}}
%%%%%%%%%%%%%%%%%%%%%%%%%%%%%%%%%%%%%%%%%%%%%%%%%%%

\newcommand{\differntllmaccuracyTable}{{
\begin{table}[t]
\centering
\captionof{table}{Perplexity ($\downarrow$) of \sysname, \sysname \texttt{(W)}, and baselines on WikiText-2 and the average accuracy ($\uparrow$) of the six common sense reasoning datasets of four different LLMs -- OPT-6.7B, LLaMA 2-7B, Mistral-7B, and Vicuna-7B -- under 20\% compression ratio. The relative performance gain compared to the best-performing baseline is marked in green color inside bracket.}
\vspace{-2mm}
\resizebox{1\textwidth}{!}{%
\begin{tabular}{c|cc|cc|cc|cc}
% \toprule[0.01pt]
\midrule
& \multicolumn{2}{c}{OPT-6.7B}   & \multicolumn{2}{c}{\textsc{LLaMA 2-7B}}    & \multicolumn{2}{c}{\textsc{Mistral-7B}}        & \multicolumn{2}{c}{\textsc{Vicuna-7B}} \\\midrule
\textsc{Method} & Perplexity\textcolor{mygreen}{$\downarrow$}    & Accuracy\textcolor{mygreen}{$\uparrow$}    & Perplexity\textcolor{mygreen}{$\downarrow$}    & Accuracy\textcolor{mygreen}{$\uparrow$}  & Perplexity\textcolor{mygreen}{$\downarrow$}    & Accuracy\textcolor{mygreen}{$\uparrow$}  & Perplexity\textcolor{mygreen}{$\downarrow$}    & Accuracy\textcolor{mygreen}{$\uparrow$} \\ \midrule
\color[HTML]{9B9B9B}Original       & \color[HTML]{9B9B9B}10.86          & \color[HTML]{9B9B9B}0.52         & \color[HTML]{9B9B9B}5.47            & \color[HTML]{9B9B9B}0.57      & \color[HTML]{9B9B9B}5.25          & \color[HTML]{9B9B9B}0.61         & \color[HTML]{9B9B9B}6.78            & \color[HTML]{9B9B9B}0.56    \\\midrule
SVD       & 66275      & 0.03         & 18192             & 0.09      & 159627        & 0.03         & 18644            & 0.05    \\
FWSVD     & 14559      & 0.06         & 2360              & 0.12      & 6357         & 0.08          & 2758              & 0.09 \\
ASVD      & 82.00      & 0.32         & 10.10             & 0.36      & 13.72         & 0.32         & 16.23             & 0.33     \\ \midrule
\xin{\sysname \texttt{(W)}}  & 16.04 (\textcolor{mygreen}{$\downarrow$80\%})         & 0.41 (\textcolor{mygreen}{$\uparrow$28\%})           &  8.50 (\textcolor{mygreen}{$\downarrow$16\%})          & 0.53 (\textcolor{mygreen}{$\uparrow$47\%})      & 10.21 (\textcolor{mygreen}{$\downarrow$26\%})         & 0.42 (\textcolor{mygreen}{$\uparrow$24\%})           &  8.41 (\textcolor{mygreen}{$\downarrow$48\%})          & 0.51 (\textcolor{mygreen}{$\uparrow$55\%})   \\
\sysname  & \textbf{14.47} (\textcolor{mygreen}{$\downarrow$82\%})         & \textbf{0.49} (\textcolor{mygreen}{$\uparrow$53\%})           &  \textbf{7.73} (\textcolor{mygreen}{$\downarrow$23\%})          & \textbf{0.54} (\textcolor{mygreen}{$\uparrow$50\%})      & \textbf{7.47} (\textcolor{mygreen}{$\downarrow$45\%})         & \textbf{0.55} (\textcolor{mygreen}{$\uparrow$72\%})           &  \textbf{7.43} (\textcolor{mygreen}{$\downarrow$54\%})          & \textbf{0.54} (\textcolor{mygreen}{$\uparrow$64\%})   \\ \midrule
% \bottomrule[0.01pt]
\end{tabular}
}
\vspace{-4mm}
\label{tab:different_llm_acc}
\end{table}
}}

%%%%%%%%%%%%%%%%%%%%%%%%%%%%%%%%%%%%%%%%%%%%%%%%%%%%%%

\newcommand{\largellmaccuracyTable}{{
\centering
\captionof{table}{Perplexity ($\downarrow$) of \sysname, \sysname \texttt{(W)}, and baselines on WikiText-2 and the average accuracy ($\uparrow$) of the six classification datasets of LLaMA-13B and LLaMA-30B under 20\% compression ratio. The relative performance gain compared to the best-performing baseline is marked in green color inside bracket.}
\vspace{-2mm}
\resizebox{0.48\textwidth}{!}{%
\begin{tabular}{c|cc|cc}
\midrule
& \multicolumn{2}{c}{\textsc{LLaMA-13B}}    & \multicolumn{2}{c}{\textsc{LLaMA-30B}}         \\\midrule
\textsc{Method}     & Perplexity\textcolor{mygreen}{$\downarrow$}    & Accuracy\textcolor{mygreen}{$\uparrow$}       & Perplexity\textcolor{mygreen}{$\downarrow$}    & Accuracy\textcolor{mygreen}{$\uparrow$}    \\ \midrule
\color[HTML]{9B9B9B}Original                 & \color[HTML]{9B9B9B}5.09        & \color[HTML]{9B9B9B}0.59           & \color[HTML]{9B9B9B}4.10         & \color[HTML]{9B9B9B}0.61 \\\midrule
SVD                 & 946.31        & 0.21           & 54.11         & 0.33            \\
FWSVD               & 15.98         & 0.43           & 20.54         & 0.42             \\
ASVD                & 6.74          & 0.54           & 22.71         & 0.44             \\ \midrule
\xin{\sysname \texttt{(W)}}            &  6.61 (\textcolor{mygreen}{$\downarrow$2\%})         & 0.54 (\textcolor{mygreen}{$\uparrow$0\%})       & 5.63 (\textcolor{mygreen}{$\downarrow$73\%})             & 0.57 (\textcolor{mygreen}{$\uparrow$30\%})          \\
\sysname            &  \textbf{6.43} (\textcolor{mygreen}{$\downarrow$5\%})         & \textbf{0.55} (\textcolor{mygreen}{$\uparrow$2\%})       &\textbf{5.14} (\textcolor{mygreen}{$\downarrow$75\%})             & \textbf{0.59} (\textcolor{mygreen}{$\uparrow$34\%})          \\ \midrule
\end{tabular}
}
\vspace{-5mm}
\label{tab:large_llm_acc}
}}

%%%%%%%%%%%%%%%%%%%%%%%%%%%%%%%%%%%%%%%%%%%%

\newcommand{\loraTable}{{
\centering
\captionof{table}{Perplexity of LLaMA-7B compressed by \sysname under 20\% to 80\% compression ratio on WikiText-2 with different updating orders.}
\vspace{1mm}
\resizebox{1\textwidth}{!}{%
\begin{tabular}{c|cccc}
\midrule
\textsc{Updating Order}     & 20\%      & 40\%      & 60\%   & 80\%   \\ \midrule
$W'_u$ first, then $W'_v$         & 7.85     & 9.32        & \textbf{13.20} (\textcolor{mygreen}{$\downarrow$1\%})  & \textbf{31.67}(\textcolor{mygreen}{$\downarrow$1\%}) \\ \midrule
$W'_v$ first, then $W'_u$          & \textbf{7.73} (\textcolor{mygreen}{$\downarrow$2\%})      & \textbf{9.27} (\textcolor{mygreen}{$\downarrow$1\%})         & 15.00    & 31.79  \\ \midrule
\end{tabular}
}
\label{tab:lora_acc}
}}

%%%%%%%%%%%%%%%%%%%%%%%%%%%%%%%%%%%%%%%%%%%%
\newcommand{\sensitivityTable}{{
\centering
\captionof{table}{Perplexity ($\downarrow$) of compressed LLaMA-7B on WikiText-2 under different compression ratios. \sysname \texttt{(W)} denotes the version of \sysname with truncation-aware data whitening only; \sysname \texttt{(U)} denotes the version of \sysname with parameter update with sequential low-rank approximation only; The relative performance gain compared to ASVD is marked in green color inside bracket.}
\vspace{5mm}
\resizebox{1\textwidth}{!}{%
% \tiny
\begin{tabular}{c|c|c|c}
% \toprule[0.01pt]
\midrule
\textsc{Method} & \multicolumn{1}{c}{\textsc{20\%}}  &\multicolumn{1}{c}{\textsc{40\%}}  &\multicolumn{1}{c}{\textsc{60\%}}     \\ \midrule          
ASVD              & 11.14                 & 1407              & 57057                        \\ \midrule
\sysname \texttt{(W)}      & 7.94 (\textcolor{mygreen}{$\downarrow$29\%})            & 13.11 (\textcolor{mygreen}{$\downarrow$99\%})          & 42.30 (\textcolor{mygreen}{$\downarrow$99\%})     \\ \midrule
\sysname \texttt{(U)}      & 10.12 (\textcolor{mygreen}{$\downarrow$9\%})            & 19.28 (\textcolor{mygreen}{$\downarrow$99\%})          & 49.88 (\textcolor{mygreen}{$\downarrow$99\%})     \\ \midrule
\sysname  & \textbf{7.73} (\textcolor{mygreen}{$\downarrow$31\%})      & \textbf{9.27} (\textcolor{mygreen}{$\downarrow$99\%})         & \textbf{15.00} (\textcolor{mygreen}{$\downarrow$99\%})           \\ \midrule
% \bottomrule[0.01pt]
\end{tabular}
}
\vspace{-4mm}
\label{tab:sensitivity}
}}

%%%%%%%%%%%%%%%%%%%%%%%%%%%%%%%%%%%%%%%%%%%%

\newcommand{\svdquantTable}{{
\centering
\captionof{table}{Perplexity ($\downarrow$) of LLaMA-7B compressed by 1-bit quantization methods and \sysname on WikiText-2. The relative performance gain compared to the best-performing baseline is marked in green.}
% \vspace{2.5mm}
\resizebox{1\textwidth}{!}{%
\begin{tabular}{c|c|c|c}
\midrule
\textsc{Method}         & \textsc{Type}         &\textsc{Memory} & \textsc{Perplexity}\\ \midrule
PB-LLM                  & Post-training         & 1.9 GB            & 104.83 \\
BiLLM                   & Post-training         & 1.5 GB            & 47.67           \\\midrule
\sysname                & Post-training         & 1.5 GB            & \textbf{47.21} (\textcolor{mygreen}{$\downarrow$1\%})        \\\midrule
OneBit                  & Training-required     & 1.3 GB            & 10.20           \\\midrule
\sysname (2-bit)       & Post-training         & 1.3 GB            & \textbf{9.83} (\textcolor{mygreen}{$\downarrow$4\%})          \\\midrule
\end{tabular}
}
\label{tab:svd_quant}
}}

%%%%%%%%%%%%%%%%%%%%%%%%%%%%%%%%%%%%%%%%%%%%%%%%%%%%%%%%%%%%%%%%%%%%%

\newcommand{\svdpruneTable}{{
\centering
\captionof{table}{Perplexity ($\downarrow$) of LLaMA-7B compressed by structured pruning methods and \sysname under various memory budget on WikiText-2. The relative performance gain compared to the best-performing baseline is marked in green.}
% \vspace{-2mm}
\resizebox{1\textwidth}{!}{%
\begin{tabular}{c|c|c|c|c}
\midrule
& \multicolumn{4}{c}{\textsc{Perplexity under various memory budget}}\\ \midrule
\textsc{Method}       & 10 GB       & 9 GB      & 8 GB    & 7 GB \\ \midrule
LLM-Pruner            & 9.88        & 12.21      & 18.94     & 21.68 \\
SliceGPT              & 8.78        & 12.73      & 16.39     & 27.41           \\
BlockPruner           & 9.40         & 12.76      & 19.78     & 43.05        \\\midrule
\sysname              & \textbf{7.92} (\textcolor{mygreen}{$\downarrow$10\%})        & \textbf{8.18} (\textcolor{mygreen}{$\downarrow$33\%})      & \textbf{8.33} (\textcolor{mygreen}{$\downarrow$49\%})     & \textbf{9.63} (\textcolor{mygreen}{$\downarrow$56\%})        \\\midrule
\end{tabular}
}
\label{tab:svd_prune}
}}

%%%%%%%%%%%%%%%%%%%%%%%%%%%%%%%%%%%%%%%%%%%%%%%%%%%%%%%%%%%%%%%%%%%%%

\newcommand{\timeTable}{{
\begin{table}[ht]
\centering
\captionof{table}{Compression time of \sysname and ASVD on LLaMA-7B under 20\% compression ratio. The relative speedup is marked in green color inside bracket.}
\vspace{-2mm}
\resizebox{0.9\textwidth}{!}{%
\begin{tabular}{ccc|ccc}
\midrule
 \multicolumn{3}{c|}{\sysname} & \multicolumn{3}{c}{ASVD}       \\ \cmidrule{1-6} 
                         \makecell{Truncation-Aware \\Data Whitening}  & \makecell{Parameter Update with Sequential \\ Low-rank Approximation}  & \textbf{Total} & Normalize & Search & \textbf{Total} \\ \midrule
                   10min      & 3.5h    & \textbf{3.5h} (\textcolor{mygreen}{$\downarrow$36\%}) & 5min          & 5.5h   & \textbf{5.5h}  \\ \midrule
\end{tabular}
}
% \vspace{-4mm}
\label{tab:time}
\end{table}
}}
%%%%%%%%%%%%%%%%%%%%%%%%%%%%%%%%%%%%%%%%%%%%%%%%%%%%%%
\newcommand{\datasetTable}{{
\centering
\captionof{table}{Performance of LLaMA-7B compressed by \sysname under 20\% compression ratio using calibration data sampled from WikiText-2 (by default in our paper) and C4 datasets. The performance on WikiText-2 and C4 are reported by perplexity ($\downarrow$), while the performance on six downstream datasets are reported by average accuracy ($\uparrow$). The performance on TruthfulQA and GSM8K are reported by BLEU score($\uparrow$) and Exact Match Accuracy ($\uparrow$) respectively. 
The relative performance gain for data sampled from one dataset compared to another is marked in green color inside bracket.
}
\vspace{-2.5mm}
\resizebox{1\textwidth}{!}{%
% \tiny
\begin{tabular}{cc|c|cc}
\midrule
WikiText-2\textcolor{mygreen}{$\downarrow$} & C4\textcolor{mygreen}{$\downarrow$} & Average\textcolor{mygreen}{$\uparrow$}  & TruthfulQA\textcolor{mygreen}{$\uparrow$}  & GSM8K\textcolor{mygreen}{$\uparrow$}\\ \midrule
\multicolumn{5}{c}{Calibration data sampled from WikiText-2 } \\ \midrule
 \textbf{7.73} (\textcolor{mygreen}{$\downarrow$1\%})       & 12.23      & \textbf{0.55} (\textcolor{mygreen}{$\uparrow$2\%})     & 0.28   & 0.08               \\\midrule
\multicolumn{5}{c}{Calibration data sampled from C4} \\ \midrule
     7.79      & \textbf{11.97} (\textcolor{mygreen}{$\downarrow$1\%})      & 0.54     & 0.28   & 0.08               \\\midrule
\end{tabular}
}
\label{tab:different_ds_acc}
}}
%%%%%%%%%%%%%%%%%%%%%%%%%%%%%%%%%%%%%%%%%%%%

\newcommand{\scalingTable}{{
\begin{table*}[t]
\centering
\vspace{6mm}
\captionof{table}{Comparison of LLaMA-3B (compressed from LLaMA-7B by \sysname) and original StableLM-3B~\citep{StableLM-3B-4E1T} trained from scratch. Both the throughput and the peak memory footprint during the inference are measured with batch size=32, sequence length = 128 on single A100 GPU. }
% \vspace{-2mm}
\resizebox{1\textwidth}{!}{%}
\begin{tabular}{c|c|c|cccccc|c|c|c}
\midrule
\textsc{Model}      & Throughput        & Peak Mem. & Openb.  & Arc\_e  & WinoG. & HellaS.   & PIQA      & MathQA    & \textbf{Average\textcolor{mygreen}{$\uparrow$}}  & TruthfulQA\textcolor{mygreen}{$\uparrow$} & GSM8K\textcolor{mygreen}{$\uparrow$} \\ \midrule
StableLM-3B         & 8463 Tokens/sec             &  9.41 GB     & 0.19   & 0.51   & 0.55    & 0.47      & 0.69      & 0.21      & 0.44  & 0.22  & 0.02                  \\ \midrule
LLaMA-3B            & 9254 Tokens/sec             &  7.43 GB     & \textbf{0.27}      & \textbf{0.54}    & \textbf{0.60}    & \textbf{0.49}      & 0.68      & 0.19      & \textbf{0.46} (\textcolor{mygreen}{$\uparrow$5\%})   & \textbf{0.23} (\textcolor{mygreen}{+ 0.01})   & \textbf{0.04} (\textcolor{mygreen}{+ 0.02})                  \\ \midrule
\end{tabular}
}
\label{tab:scaling}
% \vspace{-4mm}
\end{table*}
}}

%%%%%%%%%%%%%%%%%%%%%%%%%%%%%%%%%%%%%%%%%%%%


\newcommand{\exampleLoraTable}{{
\begin{table}[ht]
    \centering
    \vspace{-3mm}
    \caption{An example of contents generated by the fine-tuned LLaMA-7B with LoRA after being compressed by \sysname and ASVD under different compression ratios. The input is marked in bold and the normal texts are the generated sentences.} \label{tab:visualization_lora}
    \resizebox{\linewidth}{!}{
    \tiny
    \begin{tabular}{p{0.8cm}|p{5cm}|p{4.8cm}}
        \midrule
        \textsc{Ratio} & \textsc{\sysname} & \textsc{ASVD} \\
        \midrule
        Original & \multicolumn{2}{p{10.2cm}}{\textbf{What is the responsibility of the AI assistant?} The responsibility of an AI assistant is to complete the job without any hesitation. Also, the AI assistant has no problem with taking a decision and executing it. An AI assistant is responsible to complete the task and it needs to learn from his past experience to take the best possible decision without any hesitation...} \\
        \midrule
         20\% & \textbf{What is the responsibility of the AI assistant?} The AI assistant's role and responsibilities are often ambiguous and can vary depending on the context and the type of assistant being used. Some AI assistants may have preset responses to common queries, but their behavior can also depend on the user's interaction and expectations...  & \textbf{What is the responsibility of the AI assistant?} 1. To act in the same way as the programmer intended. 2. To make a value assessment of the user's intentions and offer an appropriate response. 3. To help the user carry out specific tasks or follow a strict script. 4. To make the system as independent of the user as possible...\\
        \midrule
        40\% & \textbf{What is the responsibility of the AI assistant?} 1. The AI assistant is responsible for performing tasks that can only be done by AI technology such as scheduling appointments, setting reminders and making suggestions, while also providing quick and accurate responses. 2. The AI assistant is also responsible for learning and updating its knowledge base through continuous improvement, thus making it more efficient and effective. 3. The AI assistant is responsible for handling and managing user-generated and customer-oriented input... 
        & \textbf{What is the responsibility of the AI assistant?}  1. Define the definition of the job specification. 2. Define the role of the AI assistant for the job. 3. Define the skills necessary in order to carryout the tasks assigned to the job to do the work of the same capacity in the market. 4. Determine the evaluation of the job competence of the AI assistant to ensure their professional integrity. 5. What type of AI assistant jobs are available? 6. What industries require technology assistance work to complete the tasks that it own personnel cannot perform...\\
        \midrule
        60\% & \textbf{What is the responsibility of the AI assistant?} In a nutshell, the responsibility of the AI assistant could vary depending on the task, but generally, the focus would be on automatic tasks, without the need for human intervention. Some common tasks could include setting reminders, scheduling appointments, and making routine household or productivity tasks. The AI assistant also serves as a backup or a relief system, taking on responsibilities when the user is not available ... 
        & \textbf{What is the responsibility of the AI assistant?}  2.3. ?? Brush for a discussion I wonder is it worth doing is important.2- It isn't useful just reducing labor costs; it helps employees feel a sense of connected to your attention which leads to better workplace values among staffers and leads to long relationships among org...\\
        \midrule
        80\% & \textbf{What is the responsibility of the AI assistant?}  Our Design is based on our understanding of the world, and we are actively learning, adapting and adapting, so we're always evolving new ideas, which we see to be most unique and relevant in our community... 
        & \textbf{What is the responsibility of the AI assistant?} ygua AIeltemperaturen/2, (64mbz/.3/.1/, 7.kbld.org.0/2/ In these  puthebout les bnvols n merginels ...\\
        \midrule
        % \bottomrule
    \end{tabular}
    }
\end{table}
}}
%%%%%%%%%%%%%%%%%%%%%%%%%%%%%%%%%%%%%%%%%%%%

\newcommand{\flapTable}{{
\begin{table}[t]
\centering
\captionof{table}{Perplexity ($\downarrow$) of \sysname and FLAP on WikiText-2 to compress LLaMA-7B under different compression ratios. The better performance is marked in bold. The relative performance gain of \sysname compared to FLAP is marked in green inside bracket.}
\vspace{-2mm}
\resizebox{0.8\textwidth}{!}{%
\begin{tabular}{c|c|c|c|c}
% \toprule[0.01pt]
\midrule
\textsc{Ratio (Mem.)} & \textsc{20\% (10.2GB)}   & \textsc{40\% (7.76GB)}    & \textsc{60\% (5.35GB)}        & \textsc{80\% (2.58GB)} \\\midrule
FLAP      & 7.99              & 14.43             & 106.87      & 15023         \\ \midrule
\sysname  & \textbf{7.73} (\textcolor{mygreen}{$\downarrow$3\%})    &  \textbf{9.27} (\textcolor{mygreen}{$\downarrow$36\%})               & \textbf{15.00} (\textcolor{mygreen}{$\downarrow$86\%})            &  \textbf{31.79} (\textcolor{mygreen}{$\downarrow$99\%})             \\ \midrule
% \bottomrule[0.01pt]
\end{tabular}
}
% \vspace{-5mm}
\label{tab:flap}
\end{table}
}}

%%%%%%%%%%%%%%%%%%%%%%%%%%%%%%%%%%%%%%%%%%%%

\newcommand{\morellmaccuracyTable}{{
\centering
\captionof{table}{Perplexity ($\downarrow$) on WikiText-2 of LLaMA 2-70B and LLaMA 3-8B under 20\% compression ratio.}
\vspace{-2mm}
\resizebox{0.4\textwidth}{!}{%
\begin{tabular}{c|c|c}
\midrule
\textsc{Model} & \textsc{LLaMA 2-70B}    & \textsc{LLaMA 3-8B}        \\\midrule
\color[HTML]{9B9B9B}Original                 & \color[HTML]{9B9B9B}3.32        & \color[HTML]{9B9B9B}6.14            \\\midrule
SVD                 & 19.82        & 29871            \\
FWSVD               & OOM          & 4782             \\
ASVD                & OOM          & 17.55             \\ \midrule
\sysname            &  \textbf{4.21} (\textcolor{mygreen}{$\downarrow$66\%})         & \textbf{8.16} (\textcolor{mygreen}{$\downarrow$54\%})               \\ \midrule
\end{tabular}
}
\vspace{-5mm}
\label{tab:morellm_acc}
}}

%%%%%%%%%%%%%%%%%%%%%%%%%%%%%%%%%%%%%%%%%%%%

\newcommand{\loraaddTable}{{
\begin{table}[t]
\centering
\captionof{table}{Perplexity ($\downarrow$) of \sysname with original LoRA fine-tuning (denoted as \sysname (SFT)), ASVD with sequential LoRA fine-tuning (denoted as ASVD (SFT)), and  \sysname with sequential LoRA fine-tuning (denoted as \sysname (SFT)) on WikiText-2 to compress LLaMA-7B under different compression ratios.}
\vspace{-2mm}
\resizebox{0.8\textwidth}{!}{%
\begin{tabular}{c|c|c|c|c}
% \toprule[0.01pt]
\midrule
\textsc{Ratio (Mem.)} & \textsc{20\% (10.2GB)}   & \textsc{40\% (7.76GB)}    & \textsc{60\% (5.35GB)}        & \textsc{80\% (2.58GB)} \\\midrule
\sysname (NFT)  & 7.87            & 11.98             & 16.30      & 80.23         \\ 
ASVD (SFT)  & 8.37            & 14.86             & 44.81      & 271         \\ \midrule
\sysname (SFT)  & \textbf{7.73}    &  \textbf{9.27}               & \textbf{15.00}          &  \textbf{31.79}           \\ \midrule
% \bottomrule[0.01pt]
\end{tabular}
}
\vspace{-5mm}
\label{tab:lora_add}
\end{table}
}}


%%%%%%%%%%%%%%%%%%%%%%%%%%%%%%%%%%%%%%%%%%%%

\newcommand{\orthoTable}{{
\centering
\captionof{table}{MSE between $S_{pre}S_{pre}^T$ pre-computed on WikiText-2 calibration data and $S_{fly}S_{fly}^T$ computed on-the-fly on three randomly selected weight matrices in LLaMA-7B.}
\vspace{-2mm}
\resizebox{0.4\textwidth}{!}{%
\begin{tabular}{c|c|c|c}
% \toprule[0.01pt]
\midrule
\textsc{Data} & \textsc{$W_1$} & \textsc{$W_2$}$ & \textsc{$W_3$}\\\midrule
\textsc{MSE} & 0.003 & 0.0011 & 0.0009   \\\midrule
% \bottomrule[0.01pt]
\end{tabular}
}
\vspace{-5mm}
\label{tab:orthogonality}
}}


%%%%%%%%%%%%%%%%%%%%%%%%%%%%%%%%%%%%%%%%%%%%

\newcommand{\dronelossTable}{{
\begin{table}[t]
\centering
\captionof{table}{Compression loss of the randomly generated weight and activation matrices with different shapes under 50\% compression ratio using \sysname, Drone, and the theoretical minimum.}
\vspace{-2mm}
\resizebox{0.9\textwidth}{!}{%
\begin{tabular}{c|c|c|c}
% \toprule[0.01pt]
\midrule
\textsc{Loss} & $[128\times 128] \times [128\times128]$ & $[2048\times 2048] \times [2048\times 2048]$ & $[4096\times 4096] \times [4096\times4096]$ \\\midrule
\textsc{Minimum}    &   276.1130          &     17784.2637      &   50321.9141              \\\midrule
\textsc{Drone}      &   276.1130          &     17785.6992     &   50337.2148       \\\midrule
\sysname            &   276.1130          &     17784.2676      &   50321.9727       \\\midrule
% \bottomrule[0.01pt]
\end{tabular}
}
\vspace{-4mm}
\label{tab:drone_loss}
\end{table}
}}

%%%%%%%%%%%%%%%%%%%%%%%%%%%%%%%%%%%%%%%%%%%%

\newcommand{\dronetimeTable}{{
\begin{table}[t]
\centering
\captionof{table}{Compression Time of the randomly generated weight and activation matrices with different shapes using \sysname and Drone. The compression time is measured for 5 times' compression.}
\vspace{-2mm}
\resizebox{0.9\textwidth}{!}{%
\begin{tabular}{c|c|c|c}
% \toprule[0.01pt]
\midrule
\textsc{Time} & $[128\times 128] \times [128\times128]$ & $[2048\times 2048] \times [2048\times 2048]$ & $[4096\times 4096] \times [4096\times4096]$
\\\midrule
\textsc{Drone}  &       0.07 seconds          &   5.81 seconds        &       30.35 seconds          \\\midrule
\sysname        &       0.02 seconds          &   1.98 seconds        &       10.37 seconds          \\\midrule
% \bottomrule[0.01pt]
\end{tabular}
}
\vspace{-5mm}
\label{tab:drone_time}
\end{table}
}}
% Recommended, but optional, packages for figures and better typesetting:
\usepackage{microtype}
\usepackage{graphicx}
\usepackage{subfigure}
\usepackage{caption}
\usepackage{xspace}
\usepackage{pifont}
\usepackage{makecell}
\usepackage{bm}

\usepackage{subcaption}
\usepackage{booktabs} % for professional tables
\usepackage{wrapfig}

% hyperref makes hyperlinks in the resulting PDF.
% If your build breaks (sometimes temporarily if a hyperlink spans a page)
% please comment out the following usepackage line and replace


% Attempt to make hyperref and algorithmic work together better:
\newcommand{\theHalgorithm}{\arabic{algorithm}}
\PassOptionsToPackage{numbers, compress}{natbib}
% Use the following line for the initial blind version submitted for review:
% ready for submission
\usepackage{iclr2025_conference,times}

% Optional math commands from https://github.com/goodfeli/dlbook_notation.
%%%%% NEW MATH DEFINITIONS %%%%%

\usepackage{amsmath,amsfonts,bm}

% Mark sections of captions for referring to divisions of figures
\newcommand{\figleft}{{\em (Left)}}
\newcommand{\figcenter}{{\em (Center)}}
\newcommand{\figright}{{\em (Right)}}
\newcommand{\figtop}{{\em (Top)}}
\newcommand{\figbottom}{{\em (Bottom)}}
\newcommand{\captiona}{{\em (a)}}
\newcommand{\captionb}{{\em (b)}}
\newcommand{\captionc}{{\em (c)}}
\newcommand{\captiond}{{\em (d)}}

% Highlight a newly defined term
\newcommand{\newterm}[1]{{\bf #1}}


% Figure reference, lower-case.
\def\figref#1{figure~\ref{#1}}
% Figure reference, capital. For start of sentence
\def\Figref#1{Figure~\ref{#1}}
\def\twofigref#1#2{figures \ref{#1} and \ref{#2}}
\def\quadfigref#1#2#3#4{figures \ref{#1}, \ref{#2}, \ref{#3} and \ref{#4}}
% Section reference, lower-case.
\def\secref#1{section~\ref{#1}}
% Section reference, capital.
\def\Secref#1{Section~\ref{#1}}
% Reference to two sections.
\def\twosecrefs#1#2{sections \ref{#1} and \ref{#2}}
% Reference to three sections.
\def\secrefs#1#2#3{sections \ref{#1}, \ref{#2} and \ref{#3}}
% Reference to an equation, lower-case.
\def\eqref#1{equation~\ref{#1}}
% Reference to an equation, upper case
\def\Eqref#1{Equation~\ref{#1}}
% A raw reference to an equation---avoid using if possible
\def\plaineqref#1{\ref{#1}}
% Reference to a chapter, lower-case.
\def\chapref#1{chapter~\ref{#1}}
% Reference to an equation, upper case.
\def\Chapref#1{Chapter~\ref{#1}}
% Reference to a range of chapters
\def\rangechapref#1#2{chapters\ref{#1}--\ref{#2}}
% Reference to an algorithm, lower-case.
\def\algref#1{algorithm~\ref{#1}}
% Reference to an algorithm, upper case.
\def\Algref#1{Algorithm~\ref{#1}}
\def\twoalgref#1#2{algorithms \ref{#1} and \ref{#2}}
\def\Twoalgref#1#2{Algorithms \ref{#1} and \ref{#2}}
% Reference to a part, lower case
\def\partref#1{part~\ref{#1}}
% Reference to a part, upper case
\def\Partref#1{Part~\ref{#1}}
\def\twopartref#1#2{parts \ref{#1} and \ref{#2}}

\def\ceil#1{\lceil #1 \rceil}
\def\floor#1{\lfloor #1 \rfloor}
\def\1{\bm{1}}
\newcommand{\train}{\mathcal{D}}
\newcommand{\valid}{\mathcal{D_{\mathrm{valid}}}}
\newcommand{\test}{\mathcal{D_{\mathrm{test}}}}

\def\eps{{\epsilon}}


% Random variables
\def\reta{{\textnormal{$\eta$}}}
\def\ra{{\textnormal{a}}}
\def\rb{{\textnormal{b}}}
\def\rc{{\textnormal{c}}}
\def\rd{{\textnormal{d}}}
\def\re{{\textnormal{e}}}
\def\rf{{\textnormal{f}}}
\def\rg{{\textnormal{g}}}
\def\rh{{\textnormal{h}}}
\def\ri{{\textnormal{i}}}
\def\rj{{\textnormal{j}}}
\def\rk{{\textnormal{k}}}
\def\rl{{\textnormal{l}}}
% rm is already a command, just don't name any random variables m
\def\rn{{\textnormal{n}}}
\def\ro{{\textnormal{o}}}
\def\rp{{\textnormal{p}}}
\def\rq{{\textnormal{q}}}
\def\rr{{\textnormal{r}}}
\def\rs{{\textnormal{s}}}
\def\rt{{\textnormal{t}}}
\def\ru{{\textnormal{u}}}
\def\rv{{\textnormal{v}}}
\def\rw{{\textnormal{w}}}
\def\rx{{\textnormal{x}}}
\def\ry{{\textnormal{y}}}
\def\rz{{\textnormal{z}}}

% Random vectors
\def\rvepsilon{{\mathbf{\epsilon}}}
\def\rvtheta{{\mathbf{\theta}}}
\def\rva{{\mathbf{a}}}
\def\rvb{{\mathbf{b}}}
\def\rvc{{\mathbf{c}}}
\def\rvd{{\mathbf{d}}}
\def\rve{{\mathbf{e}}}
\def\rvf{{\mathbf{f}}}
\def\rvg{{\mathbf{g}}}
\def\rvh{{\mathbf{h}}}
\def\rvu{{\mathbf{i}}}
\def\rvj{{\mathbf{j}}}
\def\rvk{{\mathbf{k}}}
\def\rvl{{\mathbf{l}}}
\def\rvm{{\mathbf{m}}}
\def\rvn{{\mathbf{n}}}
\def\rvo{{\mathbf{o}}}
\def\rvp{{\mathbf{p}}}
\def\rvq{{\mathbf{q}}}
\def\rvr{{\mathbf{r}}}
\def\rvs{{\mathbf{s}}}
\def\rvt{{\mathbf{t}}}
\def\rvu{{\mathbf{u}}}
\def\rvv{{\mathbf{v}}}
\def\rvw{{\mathbf{w}}}
\def\rvx{{\mathbf{x}}}
\def\rvy{{\mathbf{y}}}
\def\rvz{{\mathbf{z}}}

% Elements of random vectors
\def\erva{{\textnormal{a}}}
\def\ervb{{\textnormal{b}}}
\def\ervc{{\textnormal{c}}}
\def\ervd{{\textnormal{d}}}
\def\erve{{\textnormal{e}}}
\def\ervf{{\textnormal{f}}}
\def\ervg{{\textnormal{g}}}
\def\ervh{{\textnormal{h}}}
\def\ervi{{\textnormal{i}}}
\def\ervj{{\textnormal{j}}}
\def\ervk{{\textnormal{k}}}
\def\ervl{{\textnormal{l}}}
\def\ervm{{\textnormal{m}}}
\def\ervn{{\textnormal{n}}}
\def\ervo{{\textnormal{o}}}
\def\ervp{{\textnormal{p}}}
\def\ervq{{\textnormal{q}}}
\def\ervr{{\textnormal{r}}}
\def\ervs{{\textnormal{s}}}
\def\ervt{{\textnormal{t}}}
\def\ervu{{\textnormal{u}}}
\def\ervv{{\textnormal{v}}}
\def\ervw{{\textnormal{w}}}
\def\ervx{{\textnormal{x}}}
\def\ervy{{\textnormal{y}}}
\def\ervz{{\textnormal{z}}}

% Random matrices
\def\rmA{{\mathbf{A}}}
\def\rmB{{\mathbf{B}}}
\def\rmC{{\mathbf{C}}}
\def\rmD{{\mathbf{D}}}
\def\rmE{{\mathbf{E}}}
\def\rmF{{\mathbf{F}}}
\def\rmG{{\mathbf{G}}}
\def\rmH{{\mathbf{H}}}
\def\rmI{{\mathbf{I}}}
\def\rmJ{{\mathbf{J}}}
\def\rmK{{\mathbf{K}}}
\def\rmL{{\mathbf{L}}}
\def\rmM{{\mathbf{M}}}
\def\rmN{{\mathbf{N}}}
\def\rmO{{\mathbf{O}}}
\def\rmP{{\mathbf{P}}}
\def\rmQ{{\mathbf{Q}}}
\def\rmR{{\mathbf{R}}}
\def\rmS{{\mathbf{S}}}
\def\rmT{{\mathbf{T}}}
\def\rmU{{\mathbf{U}}}
\def\rmV{{\mathbf{V}}}
\def\rmW{{\mathbf{W}}}
\def\rmX{{\mathbf{X}}}
\def\rmY{{\mathbf{Y}}}
\def\rmZ{{\mathbf{Z}}}

% Elements of random matrices
\def\ermA{{\textnormal{A}}}
\def\ermB{{\textnormal{B}}}
\def\ermC{{\textnormal{C}}}
\def\ermD{{\textnormal{D}}}
\def\ermE{{\textnormal{E}}}
\def\ermF{{\textnormal{F}}}
\def\ermG{{\textnormal{G}}}
\def\ermH{{\textnormal{H}}}
\def\ermI{{\textnormal{I}}}
\def\ermJ{{\textnormal{J}}}
\def\ermK{{\textnormal{K}}}
\def\ermL{{\textnormal{L}}}
\def\ermM{{\textnormal{M}}}
\def\ermN{{\textnormal{N}}}
\def\ermO{{\textnormal{O}}}
\def\ermP{{\textnormal{P}}}
\def\ermQ{{\textnormal{Q}}}
\def\ermR{{\textnormal{R}}}
\def\ermS{{\textnormal{S}}}
\def\ermT{{\textnormal{T}}}
\def\ermU{{\textnormal{U}}}
\def\ermV{{\textnormal{V}}}
\def\ermW{{\textnormal{W}}}
\def\ermX{{\textnormal{X}}}
\def\ermY{{\textnormal{Y}}}
\def\ermZ{{\textnormal{Z}}}

% Vectors
\def\vzero{{\bm{0}}}
\def\vone{{\bm{1}}}
\def\vmu{{\bm{\mu}}}
\def\vtheta{{\bm{\theta}}}
\def\va{{\bm{a}}}
\def\vb{{\bm{b}}}
\def\vc{{\bm{c}}}
\def\vd{{\bm{d}}}
\def\ve{{\bm{e}}}
\def\vf{{\bm{f}}}
\def\vg{{\bm{g}}}
\def\vh{{\bm{h}}}
\def\vi{{\bm{i}}}
\def\vj{{\bm{j}}}
\def\vk{{\bm{k}}}
\def\vl{{\bm{l}}}
\def\vm{{\bm{m}}}
\def\vn{{\bm{n}}}
\def\vo{{\bm{o}}}
\def\vp{{\bm{p}}}
\def\vq{{\bm{q}}}
\def\vr{{\bm{r}}}
\def\vs{{\bm{s}}}
\def\vt{{\bm{t}}}
\def\vu{{\bm{u}}}
\def\vv{{\bm{v}}}
\def\vw{{\bm{w}}}
\def\vx{{\bm{x}}}
\def\vy{{\bm{y}}}
\def\vz{{\bm{z}}}

% Elements of vectors
\def\evalpha{{\alpha}}
\def\evbeta{{\beta}}
\def\evepsilon{{\epsilon}}
\def\evlambda{{\lambda}}
\def\evomega{{\omega}}
\def\evmu{{\mu}}
\def\evpsi{{\psi}}
\def\evsigma{{\sigma}}
\def\evtheta{{\theta}}
\def\eva{{a}}
\def\evb{{b}}
\def\evc{{c}}
\def\evd{{d}}
\def\eve{{e}}
\def\evf{{f}}
\def\evg{{g}}
\def\evh{{h}}
\def\evi{{i}}
\def\evj{{j}}
\def\evk{{k}}
\def\evl{{l}}
\def\evm{{m}}
\def\evn{{n}}
\def\evo{{o}}
\def\evp{{p}}
\def\evq{{q}}
\def\evr{{r}}
\def\evs{{s}}
\def\evt{{t}}
\def\evu{{u}}
\def\evv{{v}}
\def\evw{{w}}
\def\evx{{x}}
\def\evy{{y}}
\def\evz{{z}}

% Matrix
\def\mA{{\bm{A}}}
\def\mB{{\bm{B}}}
\def\mC{{\bm{C}}}
\def\mD{{\bm{D}}}
\def\mE{{\bm{E}}}
\def\mF{{\bm{F}}}
\def\mG{{\bm{G}}}
\def\mH{{\bm{H}}}
\def\mI{{\bm{I}}}
\def\mJ{{\bm{J}}}
\def\mK{{\bm{K}}}
\def\mL{{\bm{L}}}
\def\mM{{\bm{M}}}
\def\mN{{\bm{N}}}
\def\mO{{\bm{O}}}
\def\mP{{\bm{P}}}
\def\mQ{{\bm{Q}}}
\def\mR{{\bm{R}}}
\def\mS{{\bm{S}}}
\def\mT{{\bm{T}}}
\def\mU{{\bm{U}}}
\def\mV{{\bm{V}}}
\def\mW{{\bm{W}}}
\def\mX{{\bm{X}}}
\def\mY{{\bm{Y}}}
\def\mZ{{\bm{Z}}}
\def\mBeta{{\bm{\beta}}}
\def\mPhi{{\bm{\Phi}}}
\def\mLambda{{\bm{\Lambda}}}
\def\mSigma{{\bm{\Sigma}}}

% Tensor
\DeclareMathAlphabet{\mathsfit}{\encodingdefault}{\sfdefault}{m}{sl}
\SetMathAlphabet{\mathsfit}{bold}{\encodingdefault}{\sfdefault}{bx}{n}
\newcommand{\tens}[1]{\bm{\mathsfit{#1}}}
\def\tA{{\tens{A}}}
\def\tB{{\tens{B}}}
\def\tC{{\tens{C}}}
\def\tD{{\tens{D}}}
\def\tE{{\tens{E}}}
\def\tF{{\tens{F}}}
\def\tG{{\tens{G}}}
\def\tH{{\tens{H}}}
\def\tI{{\tens{I}}}
\def\tJ{{\tens{J}}}
\def\tK{{\tens{K}}}
\def\tL{{\tens{L}}}
\def\tM{{\tens{M}}}
\def\tN{{\tens{N}}}
\def\tO{{\tens{O}}}
\def\tP{{\tens{P}}}
\def\tQ{{\tens{Q}}}
\def\tR{{\tens{R}}}
\def\tS{{\tens{S}}}
\def\tT{{\tens{T}}}
\def\tU{{\tens{U}}}
\def\tV{{\tens{V}}}
\def\tW{{\tens{W}}}
\def\tX{{\tens{X}}}
\def\tY{{\tens{Y}}}
\def\tZ{{\tens{Z}}}


% Graph
\def\gA{{\mathcal{A}}}
\def\gB{{\mathcal{B}}}
\def\gC{{\mathcal{C}}}
\def\gD{{\mathcal{D}}}
\def\gE{{\mathcal{E}}}
\def\gF{{\mathcal{F}}}
\def\gG{{\mathcal{G}}}
\def\gH{{\mathcal{H}}}
\def\gI{{\mathcal{I}}}
\def\gJ{{\mathcal{J}}}
\def\gK{{\mathcal{K}}}
\def\gL{{\mathcal{L}}}
\def\gM{{\mathcal{M}}}
\def\gN{{\mathcal{N}}}
\def\gO{{\mathcal{O}}}
\def\gP{{\mathcal{P}}}
\def\gQ{{\mathcal{Q}}}
\def\gR{{\mathcal{R}}}
\def\gS{{\mathcal{S}}}
\def\gT{{\mathcal{T}}}
\def\gU{{\mathcal{U}}}
\def\gV{{\mathcal{V}}}
\def\gW{{\mathcal{W}}}
\def\gX{{\mathcal{X}}}
\def\gY{{\mathcal{Y}}}
\def\gZ{{\mathcal{Z}}}

% Sets
\def\sA{{\mathbb{A}}}
\def\sB{{\mathbb{B}}}
\def\sC{{\mathbb{C}}}
\def\sD{{\mathbb{D}}}
% Don't use a set called E, because this would be the same as our symbol
% for expectation.
\def\sF{{\mathbb{F}}}
\def\sG{{\mathbb{G}}}
\def\sH{{\mathbb{H}}}
\def\sI{{\mathbb{I}}}
\def\sJ{{\mathbb{J}}}
\def\sK{{\mathbb{K}}}
\def\sL{{\mathbb{L}}}
\def\sM{{\mathbb{M}}}
\def\sN{{\mathbb{N}}}
\def\sO{{\mathbb{O}}}
\def\sP{{\mathbb{P}}}
\def\sQ{{\mathbb{Q}}}
\def\sR{{\mathbb{R}}}
\def\sS{{\mathbb{S}}}
\def\sT{{\mathbb{T}}}
\def\sU{{\mathbb{U}}}
\def\sV{{\mathbb{V}}}
\def\sW{{\mathbb{W}}}
\def\sX{{\mathbb{X}}}
\def\sY{{\mathbb{Y}}}
\def\sZ{{\mathbb{Z}}}

% Entries of a matrix
\def\emLambda{{\Lambda}}
\def\emA{{A}}
\def\emB{{B}}
\def\emC{{C}}
\def\emD{{D}}
\def\emE{{E}}
\def\emF{{F}}
\def\emG{{G}}
\def\emH{{H}}
\def\emI{{I}}
\def\emJ{{J}}
\def\emK{{K}}
\def\emL{{L}}
\def\emM{{M}}
\def\emN{{N}}
\def\emO{{O}}
\def\emP{{P}}
\def\emQ{{Q}}
\def\emR{{R}}
\def\emS{{S}}
\def\emT{{T}}
\def\emU{{U}}
\def\emV{{V}}
\def\emW{{W}}
\def\emX{{X}}
\def\emY{{Y}}
\def\emZ{{Z}}
\def\emSigma{{\Sigma}}

% entries of a tensor
% Same font as tensor, without \bm wrapper
\newcommand{\etens}[1]{\mathsfit{#1}}
\def\etLambda{{\etens{\Lambda}}}
\def\etA{{\etens{A}}}
\def\etB{{\etens{B}}}
\def\etC{{\etens{C}}}
\def\etD{{\etens{D}}}
\def\etE{{\etens{E}}}
\def\etF{{\etens{F}}}
\def\etG{{\etens{G}}}
\def\etH{{\etens{H}}}
\def\etI{{\etens{I}}}
\def\etJ{{\etens{J}}}
\def\etK{{\etens{K}}}
\def\etL{{\etens{L}}}
\def\etM{{\etens{M}}}
\def\etN{{\etens{N}}}
\def\etO{{\etens{O}}}
\def\etP{{\etens{P}}}
\def\etQ{{\etens{Q}}}
\def\etR{{\etens{R}}}
\def\etS{{\etens{S}}}
\def\etT{{\etens{T}}}
\def\etU{{\etens{U}}}
\def\etV{{\etens{V}}}
\def\etW{{\etens{W}}}
\def\etX{{\etens{X}}}
\def\etY{{\etens{Y}}}
\def\etZ{{\etens{Z}}}

% The true underlying data generating distribution
\newcommand{\pdata}{p_{\rm{data}}}
% The empirical distribution defined by the training set
\newcommand{\ptrain}{\hat{p}_{\rm{data}}}
\newcommand{\Ptrain}{\hat{P}_{\rm{data}}}
% The model distribution
\newcommand{\pmodel}{p_{\rm{model}}}
\newcommand{\Pmodel}{P_{\rm{model}}}
\newcommand{\ptildemodel}{\tilde{p}_{\rm{model}}}
% Stochastic autoencoder distributions
\newcommand{\pencode}{p_{\rm{encoder}}}
\newcommand{\pdecode}{p_{\rm{decoder}}}
\newcommand{\precons}{p_{\rm{reconstruct}}}

\newcommand{\laplace}{\mathrm{Laplace}} % Laplace distribution

\newcommand{\E}{\mathbb{E}}
\newcommand{\Ls}{\mathcal{L}}
\newcommand{\R}{\mathbb{R}}
\newcommand{\emp}{\tilde{p}}
\newcommand{\lr}{\alpha}
\newcommand{\reg}{\lambda}
\newcommand{\rect}{\mathrm{rectifier}}
\newcommand{\softmax}{\mathrm{softmax}}
\newcommand{\sigmoid}{\sigma}
\newcommand{\softplus}{\zeta}
\newcommand{\KL}{D_{\mathrm{KL}}}
\newcommand{\Var}{\mathrm{Var}}
\newcommand{\standarderror}{\mathrm{SE}}
\newcommand{\Cov}{\mathrm{Cov}}
% Wolfram Mathworld says $L^2$ is for function spaces and $\ell^2$ is for vectors
% But then they seem to use $L^2$ for vectors throughout the site, and so does
% wikipedia.
\newcommand{\normlzero}{L^0}
\newcommand{\normlone}{L^1}
\newcommand{\normltwo}{L^2}
\newcommand{\normlp}{L^p}
\newcommand{\normmax}{L^\infty}

\newcommand{\parents}{Pa} % See usage in notation.tex. Chosen to match Daphne's book.

\DeclareMathOperator*{\argmax}{arg\,max}
\DeclareMathOperator*{\argmin}{arg\,min}

\DeclareMathOperator{\sign}{sign}
\DeclareMathOperator{\Tr}{Tr}
\let\ab\allowbreak



\usepackage[utf8]{inputenc} % allow utf-8 input
\usepackage[T1]{fontenc}    % use 8-bit T1 fonts
\usepackage{hyperref}       % hyperlinks
\usepackage{url}            % simple URL typesetting
\usepackage{amsfonts}       % blackboard math symbols
\usepackage{nicefrac}       % compact symbols for 1/2, etc.
\usepackage{xcolor}         % colors


% If accepted, instead use the following line for the camera-ready submission:

% For theorems and such
\usepackage{amsmath}
\usepackage{amssymb}
\usepackage{mathtools}
\usepackage{amsthm}
\usepackage{multirow}
\usepackage{graphicx}
\usepackage{lscape}
\usepackage{algorithm}
\usepackage{algpseudocode}
% if you use cleveref..
\usepackage[capitalize,noabbrev]{cleveref}

%%%%%%%%%%%%%%%%%%%%%%%%%%%%%%%%
% THEOREMS
%%%%%%%%%%%%%%%%%%%%%%%%%%%%%%%%
\theoremstyle{plain}
\newtheorem{theorem}{Theorem}[section]
\newtheorem{proposition}[theorem]{Proposition}
\newtheorem{lemma}[theorem]{Lemma}
\newtheorem{corollary}[theorem]{Corollary}
\theoremstyle{definition}
\newtheorem{definition}[theorem]{Definition}
\newtheorem{assumption}[theorem]{Assumption}
\theoremstyle{remark}
\newtheorem{remark}[theorem]{Remark}
\newcommand{\sysname}{\texttt{SVD-LLM}\xspace}
% \newcommand{\sysnamew}{\texttt{SVDLLM-W}\xspace}
% \newcommand{\sysnamewl}{\texttt{SVDLLM-WL}\xspace}
\newcommand{\remaining}[1]{\textcolor{red}{Mark: #1}}
% Todonotes is useful during development; simply uncomment the next line
%    and comment out the line below the next line to turn off comments
%\usepackage[disable,textsize=tiny]{todonotes}
\usepackage[textsize=tiny]{todonotes}